\chapter{Collections}
\label{ch:collections}

\standfirst{One protocol, many implementations. Learn \ct{do:},
\ct{collect:}, \ct{select:}, \ct{detect:} and \ct{inject:into:} and you can
use every collection in the system, including ones nobody has written yet.}

\section{The hierarchy}

\begin{plain}
Collection
  Bag
  Set
    Dictionary
      IdentityDictionary
        MethodDictionary
  SequenceableCollection
    ArrayedCollection
      Array
      String
        Symbol
      Text
      ByteArray
      Interval
    LinkedList
    OrderedCollection
      SortedCollection
\end{plain}

The important lines are the abstract ones. \ct{Collection} defines nearly all
the protocol in terms of \ct{do:}, so a class that implements \ct{do:} and
\ct{add:} inherits about eighty working methods.
\ct{SequenceableCollection} adds everything that depends on order and
integer indices.

\begin{stnote}
The hierarchy above is 1983's, where \ct{Dictionary} is a subclass of
\ct{Set}. The \ct{pharo-collections} profile substitutes Pharo's versions, in
which \ct{HashedCollection} sits under \ct{Collection} with \ct{Set} and
\ct{Dictionary} as siblings. Both pass their own tests. Which you get depends
on the profile you built with.
\end{stnote}

\section{Choosing one}

\begin{center}
\small
\begin{tabular}{@{}lp{0.60\textwidth}@{}}
\toprule
\ctp{Array} & fixed size, integer-indexed, fastest. \ctp{Array new: 10}\\
\ctp{OrderedCollection} & grows and shrinks at either end. The default choice\\
\ctp{Set} & no duplicates, no order, membership by \ctp{=} and \ctp{hash}\\
\ctp{IdentitySet} & membership by \ctp{==}\\
\ctp{Dictionary} & key to value, keys compared by \ctp{=}\\
\ctp{IdentityDictionary} & keys compared by \ctp{==}\\
\ctp{Bag} & a Set that counts duplicates\\
\ctp{SortedCollection} & keeps itself sorted by a block\\
\ctp{Interval} & an arithmetic sequence; what \ctp{1 to: 10} answers\\
\ctp{String} & a byte-indexed collection of Characters\\
\ctp{Symbol} & an interned, immutable String\\
\bottomrule
\end{tabular}
\end{center}

\section{Making them}

\begin{st}
Array new: 5                      "5 nils"
Array with: 1 with: 2             "up to six with:s"
#(1 2 3)                          "literal; shared, do not mutate"
{ 1. 2 + 3 }                      "built fresh at run time"

OrderedCollection new
OrderedCollection withAll: #(1 2 3)
(1 to: 5) asOrderedCollection

Set new
Set withAll: #(1 1 2)             "2 elements"

Dictionary new
(Dictionary new) at: #a put: 1; yourself

SortedCollection sortBlock: [:a :b | a > b]
\end{st}

Every collection also answers \ct{asArray}, \ct{asOrderedCollection},
\ct{asBag}, \ct{asSet} and \ct{asSortedCollection}, so converting is one
message.

\begin{stwarn}
A literal array \ct|#(1 2 3)| is compiled once and lives in the method. If you
modify it, you have modified the method's constant, and every later evaluation
sees the change. Use \ct|{ }| or \ct{copy} when you intend to mutate.
\end{stwarn}

\section{The five you will use constantly}

\subsection{\texttt{do:}}

\begin{st}
#(1 2 3) do: [:each | Transcript show: each printString].
#(1 2 3) do: [:each | ... ] separatedBy: [ ... ].
#(1 2 3) reverseDo: [:each | ... ].
#(1 2 3) doWithIndex: [:each :i | ... ].
aDictionary keysAndValuesDo: [:k :v | ... ].
aDictionary associationsDo: [:assoc | ... ].
\end{st}

\subsection{\texttt{collect:} --- transform}

\begin{st}
#(1 2 3) collect: [:x | x * 10]        "(10 20 30 )"
\end{st}

Answers a collection of the same kind---collecting over an
\ct{OrderedCollection} gives an \ct{OrderedCollection}. The exceptions are the
ones where the same kind makes no sense: over an \ct{Interval} you get an
\ct{Array}, because an interval of squares is not an interval.

\subsection{\texttt{select:} and \texttt{reject:} --- filter}

\begin{st}
#(1 2 3 4) select: [:x | x even]        "(2 4 )"
#(1 2 3 4) reject: [:x | x even]        "(1 3 )"
\end{st}

\subsection{\texttt{detect:} --- find the first}

\begin{st}
#(1 2 3) detect: [:x | x > 1]                "2"
#(1 2 3) detect: [:x | x > 9] ifNone: [0]    "0"
\end{st}

\ct{detect:} without \ct{ifNone:} raises an error when nothing matches. Prefer
the two-keyword form; the block is the whole point of the design.

\subsection{\texttt{inject:into:} --- fold}

\begin{st}
(1 to: 10) inject: 0 into: [:sum :each | sum + each]      "55"
#('a' 'b') inject: '' into: [:acc :each | acc , each]     "'ab'"
\end{st}

The first argument is the starting value; the block gets the accumulator and
the next element and answers the new accumulator.

\section{The rest of the protocol}

\begin{center}
\small
\begin{tabular}{@{}ll@{}}
\toprule
\multicolumn{2}{@{}l}{\textit{asking}}\\
\ctp{size} \ctp{isEmpty} \ctp{notEmpty} & \\
\ctp{includes:} \ctp{occurrencesOf:} & \\
\ctp{anySatisfy:} \ctp{allSatisfy:} \ctp{noneSatisfy:} & \\
\ctp{count:} & how many satisfy a block\\
\ctp{isCollection} & \\
\midrule
\multicolumn{2}{@{}l}{\textit{arithmetic}}\\
\ctp{sum} \ctp{sum:} \ctp{max} \ctp{min} \ctp{average} & \\
\midrule
\multicolumn{2}{@{}l}{\textit{changing}}\\
\ctp{add:} \ctp{addAll:} \ctp{remove:} \ctp{remove:ifAbsent:} & \\
\ctp{addFirst:} \ctp{addLast:} \ctp{removeFirst} \ctp{removeLast} \ctp{removeIndex:} & ordered only\\
\midrule
\multicolumn{2}{@{}l}{\textit{sequenceable}}\\
\ctp{first} \ctp{last} \ctp{first:} \ctp{last:} \ctp{allButFirst} \ctp{allButLast} & \\
\ctp{at:} \ctp{at:put:} \ctp{at:ifAbsent:} & \\
\ctp{indexOf:} \ctp{copyFrom:to:} \ctp{reverse} \ctp{reversed} & \\
\ctp{,} & concatenate\\
\ctp{with:collect:} & zip two collections\\
\midrule
\multicolumn{2}{@{}l}{\textit{sorting and shaping}}\\
\ctp{sorted} \ctp{sorted:} \ctp{asSortedCollection:} & \\
\ctp{flatCollect:} & collect and flatten one level\\
\ctp{copyWithout:} \ctp{copyWith:} & \\
\ctp{ifEmpty:} \ctp{ifNotEmpty:} \ctp{anyOne} & \\
\bottomrule
\end{tabular}
\end{center}

\begin{stnote}[One design decision worth knowing]
Of the forty-years-of-protocol methods this system adds, \ct{sorted} and
\ct{sorted:} are the two that do not always answer the kind of collection
they were sent to, and the reason is that sometimes there is no such kind: a
sorted Set is not a Set, so an \ct{Array} is the best that can be said about
it. Send them to a sequenceable collection and the species is obvious and is
what you get---\ct{'hello' sorted} is \ct{'ehllo'} and not an \ct{Array} of
five Characters. Everything that preserves shape---\ct{select:},
\ct{reject:}, \ct{collect:}, \ct{flatCollect:}, \ct{copyWithout:}---goes on
using \ct{species}, because there the 1983 answer was already right.
\end{stnote}

\section{Dictionaries}

\begin{st}
| d |
d := Dictionary new.
d at: #one put: 1.
d at: #two put: 2.
d at: #one                       "1"
d at: #three ifAbsent: [0]       "0"
d at: #three ifAbsentPut: [3]    "3, and now it is in there"
d removeKey: #one                "answers the ASSOCIATION #one->1"
d includesKey: #two              "true"
d keys                           "a Set of the keys"
d values                         "a Bag of the values"
d associations                   "key -> value pairs"
d keysAndValuesDo: [:k :v | ... ]
d keysDo: [:k | ... ]
d associationsDo: [:a | ... ]
\end{st}

\ct{do:} on a Dictionary iterates the \emph{values}, not the keys and not the
associations. That trips everyone once.

An \ct{Association} is an ordinary object made with \ct{->}:

\begin{st}
#a -> 1                      "a->1"
(#a -> 1) key                "#a"
(#a -> 1) value              "1"
\end{st}

\section{Two rules about \texttt{=} and \texttt{hash}}

\begin{enumerate}
\item If you override \ct{=}, override \ct{hash}, and make \ct{a = b} imply
  \ct{a hash = b hash}. A \ct{Set} that disagrees with you will hold two
  objects you consider equal, and \ct{includes:} will answer \ct{false} for
  something that is in there.
\item Do not mutate an object while it is a key in a \ct{Dictionary} or an
  element of a \ct{Set}. It is filed under its old hash and will not be found
  again.
\end{enumerate}

\section{And one rule about threads}

\begin{stwarn}
\textbf{The base collections are not synchronised.} \ct{OrderedCollection},
\ct{Dictionary}, \ct{Set} and \ct{Bag} may be read by many workers at once and
written by exactly one, or you must guard them yourself. This is deliberate:
paying for a lock on every access to serve the rare shared case is the wrong
default, which is the lesson Java learned from \ct{Vector}. Use \ct{Mutex},
\ct{Monitor} or \ct{SharedQueue} when a collection really is shared---see
Chapter~\ref{ch:parallel}.
\end{stwarn}

Three classes in the library are the exception, and each says why in its own
chapter: \ct{JSONObject} and \ct{JSONArray}
(Chapter~\ref{ch:json}) each hold a \ct{Mutex}, because a JSON document is
the object a request arrives as and is then handed to whatever answers it; and
\ct{DbSchemaGraph} (Chapter~\ref{ch:database}) holds one because a schema is
read once and given to every connection. They are exceptions on the same
argument that makes the rule---the default is right, and these are the cases
where the shared use is not rare but expected.
